%
% teil3.tex -- Beispiel-File für Teil 3
%
% (c) 2020 Prof Dr Andreas Müller, Hochschule Rapperswil
%
% !TEX root = ../../buch.tex
% !TEX encoding = UTF-8
%
\section{Eulers Beweis}
\label{basel:section:beweis}
 
Für die heutige Zeit ist der Beweis von Euler nicht ganz aktuell.
Die Standards für einen korrekten Beweis waren zu seiner Zeit andere als heute.
 
\subsection{Linearfaktoren}
\label{basel:subsection:linearfaktoren}
 
Wenn ein Polynom $p(x)$ den Grad $n$ hat und die Nullstellen $a_1, a_2, a_3, \ldots, a_n$ besitzt,
dann kann es in seine Linearfaktoren zerlegt werden.
Man erhält
\[
    p(x) = b(x - a_1)(x - a_2)(x - a_3) \cdots (x - a_n),
    \quad\text{mit}\quad
    b = p(0) \Big/ \prod_{i=1}^{n} (-a_i).
\]
Teilt man durch das Produkt aller Nullstellen, ergibt sich
\[
    p(x)
    =
    p(0)
    \left(1 - \frac{x}{a_1}\right)
    \left(1 - \frac{x}{a_2}\right)
    \left(1 - \frac{x}{a_3}\right)
    \cdots
    \left(1 - \frac{x}{a_n}\right).
\]
Euler setzte voraus, dass dies bei Potenzreihen analog zu Polynomen einfach mit unendlich vielen Nullstellen funktioniert.
Ob das stimmt, hätte er beweisen müssen -- hat er aber nicht.
 
\subsection{Der Kardinalsinus}
\label{basel:subsection:kardinalsinus}
 
Der \emph{Kardinalsinus}
\[
    \operatorname{si}(x) = \frac{\sin x}{x}
\]
hat bis auf die Null die gleichen Nullstellen wie der Sinus.
Euler stellt nun den Kardinalsinus als Produkt der Linearfaktoren dar,
\begin{align*}
    \operatorname{si}(x)
    &\stackrel{?}{=}
    \prod_{k=1}^{\infty}
    \left(1 - \frac{x}{k\pi}\right)
    \left(1 + \frac{x}{k\pi}\right)
    \\
    &=
    \left(1 - \frac{x}{\pi}\right)
    \left(1 + \frac{x}{\pi}\right)
    \left(1 - \frac{x}{2\pi}\right)
    \left(1 + \frac{x}{2\pi}\right)
    \left(1 - \frac{x}{3\pi}\right)
    \left(1 + \frac{x}{3\pi}\right)
    \cdots.
\end{align*}
Das darf man im Allgemeinen nicht einfach so machen.
150 Jahre später bewies Karl Weierstrass, dass es in diesem Fall aber funktioniert.
 
Mit Hilfe der 3.\ binomischen Formel werden jeweils die Faktoren zusammengefasst,
welche sich nur durch das Vorzeichen unterscheiden.
Man erhält diese unendlichen Produkte
\[
    \left(1 - \frac{x^2}{\pi^2}\right)
    \left(1 - \frac{x^2}{4\pi^2}\right)
    \left(1 - \frac{x^2}{9\pi^2}\right)
    \cdots
    =
    \prod_{k=1}^{\infty}
    \left(1 - \frac{x^2}{k^2 \pi^2}\right).
\]
Damit ist aber eigentlich ein Grenzwert endlicher Produkte gemeint,
\[
    \lim_{n \to \infty}
    \prod_{k=1}^{n}
    \left(1 - \frac{x^2}{k^2 \pi^2}\right),
\]
welcher jetzt ausmultipliziert werden kann.
Ein Produkt mit $n$ Faktoren ergibt jeweils ein Polynom mit dem Grad $2n$,
\[
    \lim_{n \to \infty}
    \sum_{i=1}^{2n} a_{n,i}\, x^i,
    \qquad\text{mit}\quad
    a_{n,2}
    =
    -\sum_{k=1}^{n} \frac{1}{k^2 \pi^2}
    =
    -\frac{1}{\pi^2} \sum_{k=1}^{n} \frac{1}{k^2}.
\]
Der Grenzwert wäre dann also
\[
    \sum_{i=0}^{\infty} a_i\, x^i,
    \qquad\text{mit}\quad
    a_2 = -\frac{1}{\pi^2} \sum_{k=1}^{\infty} \frac{1}{k^2},
\]
und bis auf den Vorfaktor entspricht dieser dem Basel-Problem.
 
\subsection{Die Taylorreihe}
\label{basel:subsection:taylor}
 
Euler benutzte ebenfalls noch eine andere Darstellung des Kardinalsinus.
Er ging von der Taylorreihe für den Sinus aus und teilte durch $x$,
\[
    \operatorname{si}(x)
    =
    \frac{\sin x}{x}
    =
    \sum_{n=0}^{\infty} \frac{(-1)^n x^{2n}}{(2n+1)!}
    =
    1 - \frac{x^2}{3!} + \frac{x^4}{5!} - \frac{x^6}{7!} + \frac{x^8}{9!} \pm \cdots.
\]
 
\subsection{Der Beweis}
\label{basel:subsection:der-beweis}
 
Diese zwei Darstellungen des Kardinalsinus als Potenzreihen müssen natürlich identisch sein.
Darum kann man einen Koeffizientenvergleich machen.
Für $x^2$ mit $a_2$ erhält man
\[
    -\frac{1}{3!}
    =
    -\frac{1}{\pi^2} \sum_{k=1}^{\infty} \frac{1}{k^2},
\]
die zu beweisende Identität.
Jetzt muss nur noch auf beiden Seiten mit $-\pi^2$ multipliziert werden, und man erhält
\begin{equation}
    \frac{\pi^2}{6}
    =
    \sum_{k=1}^{\infty} \frac{1}{k^2}.
    \label{basel:equation:ergebnis}
\end{equation}
Somit stimmt der Beweis von Euler~\eqref{basel:equation:ergebnis}, auch wenn er nicht die heutigen Standards erfüllt.